\chapter{MongoDB Setup}

\section{Core Concepts}

\begin{itemize}
  \item \textbf{MongoDB} --- a document-oriented NoSQL database; stores data as flexible, JSON-like documents instead of rows in tables.
  \item \textbf{MongoDB Atlas} --- MongoDB's official managed cloud hosting service; provisions clusters without you managing servers.
  \item \textbf{Database} --- a logical container for collections (e.g. \texttt{taskmanager}).
  \item \textbf{Collection} --- a group of documents, analogous to a SQL table (e.g. \texttt{tasks}, \texttt{users}).
  \item \textbf{Document} --- a single BSON record inside a collection, analogous to a row (e.g. one task).
  \item \textbf{ObjectId} --- MongoDB's default 12-byte unique identifier, auto-assigned to a document's \texttt{\_id} field.
  \item \textbf{BSON} --- Binary JSON; the binary-encoded format MongoDB actually stores documents in on disk (adds types like Date and ObjectId that plain JSON lacks).
\end{itemize}

\begin{diagrambox}[MongoDB Hierarchy]
\begin{verbatim}
MongoDB Server / Cluster
        |
        v
    Database            (e.g. "taskmanager")
        |
        v
    Collection           (e.g. "tasks", "users")
        |
        v
    Document             ({ _id, title, status, ... })
        |
        v
    Fields               (title: "Fix bug", status: "todo")
\end{verbatim}
\end{diagrambox}

\section{Connecting with Mongoose}

\whatwhyhow{
Mongoose is an ODM (Object Data Modeling) library that sits between your Node.js code and MongoDB, giving you schemas, validation, and query helpers.
}{
Raw MongoDB driver calls are untyped and repetitive. Mongoose lets you define a schema once and get consistent validation and query building everywhere.
}{
Create a single connection module (\texttt{config/db.js}) that reads the connection string from environment variables and connects once at startup.
}

\subsection{config/db.js}

\begin{lstlisting}[language=JavaScript]
const mongoose = require("mongoose");

const connectDB = async () => {
  try {
    const conn = await mongoose.connect(process.env.MONGO_URI);
    console.log(`MongoDB connected: ${conn.connection.host}`);
  } catch (error) {
    console.error(`MongoDB connection error: ${error.message}`);
    process.exit(1);
  }
};

module.exports = connectDB;
\end{lstlisting}

\subsection{server.js (excerpt)}

\begin{lstlisting}[language=JavaScript]
const connectDB = require("./config/db");

connectDB().then(() => {
  app.listen(PORT, () => console.log(`Server running on port ${PORT}`));
});
\end{lstlisting}

\section{MongoDB Atlas Connection String}

An Atlas connection string follows the \texttt{mongodb+srv://} format:

\begin{lstlisting}[language=JavaScript]
// .env
MONGO_URI=mongodb+srv://<username>:<password>@cluster0.mongodb.net/taskmanager?retryWrites=true&w=majority
\end{lstlisting}

\begin{notebox}[NOTE]
Atlas requires you to whitelist IP addresses under Network Access before any connection succeeds; for local development add your current IP (or \texttt{0.0.0.0/0} temporarily for testing only).
\end{notebox}

\begin{warnbox}[WARNING]
Never commit \texttt{.env} or hardcode the connection string with real credentials into source control. Add \texttt{.env} to \texttt{.gitignore} and provide a \texttt{.env.example} with placeholder values instead.
\end{warnbox}
